\chapter{CONCLUSION AND FUTURE SCOPE}
\thispagestyle{empty}

This chapter concludes the project, lists its key achievements and
limitations, and clearly separates the deferred future work from the
implemented functionality.

\clearpage

\section{Conclusion}

The project set out to automate the pipeline from raw syllabus and study
material to a delivered, evaluated examination inside a single tenant-safe
platform. That objective is realized. The implemented CatLium EduTech
platform provides identity and tenancy with role based access, AI-assisted
and teacher-confirmed syllabus structuring, a distributed OCR pipeline with
page-level fidelity, a deterministic Material Intelligence engine, AI
generation validated on both worker and API boundaries, a question domain
with deterministic extraction, paper patterns and question papers with
coverage gating, the complete examination and attempt lifecycle, and
ungraded practice---all served to administrators, teachers, and students
through a polished web application.

The architecture decisions described in the proposal were preserved: the API
is a modular monolith, PostgreSQL is the only primary database, RabbitMQ
carries the async workload, the OCR service is the only separately deployable
component, and every data access is tenant-scoped. The system is validated by
188 API unit tests, 118 Python tests across the OCR and worker services, 15
web tests, and green end-to-end harnesses that run against the live stack.
The repository is documented so that development can be resumed from any
checkpoint.

\section{Key Achievements}

\begin{enumerate}
  \item An end-to-end, browser-usable multi-tenant platform with strict
        tenant and role enforcement.
  \item Deterministic OCR, Material Intelligence, pattern/question
        extraction, and objective-question grading, complemented by
        schema-validated, human-reviewed AI generation.
  \item Asynchronous job infrastructure that keeps the API responsive and
        supports retries and racing-safe cancellation.
  \item A continuously validated codebase with per-layer unit tests, type
        checking, linting, and live end-to-end verification after every code
        change.
\end{enumerate}

\section{Limitations}

The remaining limitations of the implemented system are: text and essay
answers are stored but not graded; practice is deliberately ungraded;
per-attempt (N-of-M) mechanics are not implemented; online answer-sheet
capture, OMR, and on-screen marking are absent; Redis is declared but unused;
and no continuous-integration pipeline or quantitative performance benchmarks
exist in the repository.

\section{Future Scope}

The following items are intentionally future work and are separated from the
implemented functionality:

\begin{enumerate}
  \item \textbf{Checking system.} Online answer-sheet capture, optical mark
        recognition, and on-screen marking, together with automatic grading
        of text and essay answers.
  \item \textbf{Practice and attempts.} Practice scoring, and per-attempt
        (N-of-M) mechanics with automatic replenishment after selection.
  \item \textbf{Question domain.} Question versioning and question-set
        delete/merge operations.
  \item \textbf{Academics.} The redesigned per-topic, per-resource academic
        export, and derived study-content generation from enhanced materials.
  \item \textbf{Infrastructure.} Adoption of the declared Redis service for
        caching, a continuous-integration pipeline, and measured performance
        and OCR-accuracy benchmarks.
\end{enumerate}